\documentclass[12px]{article}

\title{Lezione 15 Metodi Numerici di Approssimazione}
\date{2026-04-28}
\author{Federico De Sisti}

\input{../../setup.tex}

\begin{document}
	\maketitle
	\newpage
	\subsection{boh}$
	 \begin{cases}
		 \dot {y_1} = y_2\\
		 \dot {y_2} = -\sin y_1\\
		y_1(0) = y_{1,0}\\
		y_2(0) = y_{2,0}\\
	 \end{cases}$\\
$y_1$ è l'angolo e $y_2$ è la velocità angolare\\
La consegna è
\begin{enumerate}
	\item Grafico di  $E(y_1,y_2)$
	\item Calcolo della soluziona numerica $U^n = (u_1^n, u_2^n)$ e con, Eulero in avanti, Heun, RK4, Eulero all'indietro, Crank-Nicholson (per quelli impilici usare il metodo i newton)
	\item Grafico della soluzione sovrapposto a quello di $E$ IMMAGINE
	\item Grafico dell'energia discreta nel tempo confrontato con l'energia esatta IMMAGINE
\end{enumerate}
Siccome tutti i metodi che abbiamo visto perdono o guadagnano energia non ci aspettiamo che segua le curve di livello.\\
\textbf{Consiglio}\\
Strutturare il codice in moduli
\begin{itemize}
	\item Function per la dinamica $F:\R^2 \rightarrow \R^2$
	\item Un blocco di codice per ogni metodo
\end{itemize}



	
\end{document}
